% archon:covers AlgebraicJacobian/Albanese/CoheightBridge.lean

\chapter{Coheight--Krull dim bridge for scheme points}
\label{chap:Albanese_CoheightBridge}

% STRATEGY NOTE (iter-183). This chapter is a project-side scaffold sitting
% between \cref{chap:Albanese_CodimOneExtension} and the Riemann--Roch
% Weil-divisor machinery of \cref{chap:RiemannRoch_WeilDivisor}. The
% directive originates from the iter-182 plan-phase api-alignment consult
% \texttt{analogies/stacks-00tt-coheight.md} (Decision 2), which determined
% that Mathlib at the project's pinned commit (\texttt{b80f227}) ships every
% prerequisite for the scheme-level coheight \(=\) Krull-dim-of-stalk
% identification but does \emph{not} ship the bridge itself. The bridge is
% \emph{not} a published-textbook lemma; it is the project-side assembly of
% four standard Mathlib API pieces:
%   1. \texttt{Order.coheight\_orderIso} -- coheight transport across order
%      isomorphisms (Mathlib \texttt{Order.KrullDimension}).
%   2. \texttt{Topology.Inseparable.Specializes.mem\_open} -- opens contain
%      every generisation of their points (Mathlib \texttt{Topology.Inseparable}).
%   3. \texttt{IsLocalization.AtPrime.ringKrullDim\_eq\_height} -- the algebra
%      bridge between localisation Krull dim and prime-ideal height
%      (Mathlib \texttt{RingTheory.Ideal.Height}).
%   4. \texttt{IsAffineOpen.isLocalization\_stalk} -- the affine-stalk
%      \(=\) localisation identification (Mathlib \texttt{AlgebraicGeometry.AffineScheme}).
% The Lean target file \texttt{AlgebraicJacobian/Albanese/CoheightBridge.lean}
% does not yet exist; this chapter is the prover-ready specification for
% creating it. None of the four blocks below has a direct Hartshorne / Milne /
% Mumford / Stacks tag that quotes the exact statement, so the blocks are
% packaged as Archon-original assembly lemmas with explicit Mathlib API
% pointers in lieu of \texttt{\% SOURCE QUOTE:} blocks.

\section{Setup and motivation}
\label{sec:coheight_setup}

A point \(z\) of a scheme \(X\) has two natural ``dimension'' invariants attached
to it. The \emph{Krull dimension of the local ring at \(z\)}, written
\(\dim \mathcal{O}_{X,z}\), is an algebraic quantity: the supremum of the lengths
of chains of prime ideals in the stalk \(\mathcal{O}_{X,z}\). The
\emph{coheight} of \(z\) in \(X\), written \(\operatorname{coheight}_X(z)\), is a
topological quantity attached to the specialisation preorder on the underlying
topological space of \(X\): the supremum of the lengths of strict-generisation
chains \(z = a_0 \prec a_1 \prec \cdots \prec a_n\) in \(X\), where \(\prec\) is the
strict specialisation order \(x \prec y \iff y \rightsquigarrow x \wedge y \neq x\).

These two invariants agree:
\[
  \dim \mathcal{O}_{X,z} \;=\; \operatorname{coheight}_X(z),
\]
which is the headline theorem \cref{thm:ringKrullDim_stalk_eq_coheight} of this
chapter. The identification is folklore in algebraic geometry --- it is the
statement that the dimension of the local ring at \(z\) counts irreducible
closed subsets of \(X\) containing \(z\) --- but at the project's pinned Mathlib
commit (\texttt{b80f227}) the equation is not packaged as a single declaration
and must be assembled project-side from four Mathlib API pieces.

\paragraph{Why the bridge is needed.} Two downstream files in the project
currently work around the missing bridge by threading explicit
\texttt{[Ring.KrullDimLE 1 (X.presheaf.stalk \_)]} instance arguments at every
codim-\(1\) use site:
\begin{itemize}
  \item \cref{chap:Albanese_CodimOneExtension} carries an internal conjunction
    \[
      \text{hreg\_dim : IsRegularLocalRing\ } \mathcal{O}_{X,z} \;\wedge\;
      \operatorname{ringKrullDim}(\mathcal{O}_{X,z}) = 1
    \]
    inside the private helper
    \texttt{smooth\_codim\_one\_maximalIdeal\_isPrincipal\_and\_ne\_bot};
    the Krull-dim conjunct is supplied by \emph{this chapter's}
    \cref{lem:ringKrullDimLE_of_coheight_eq_one}, leaving only the
    irreducible \texttt{IsRegularLocalRing} half (which is the separate
    Stacks tag \texttt{00TT} sub-project, out of scope here).
  \item \cref{chap:RiemannRoch_WeilDivisor}'s order-at-a-prime-divisor map
    \texttt{Scheme.RationalMap.order} currently carries an explicit
    \texttt{[Ring.KrullDimLE 1 (X.presheaf.stalk Y.point)]} instance argument;
    after \cref{lem:ringKrullDimLE_of_coheight_eq_one} lands, the instance is
    synthesisable from the codim-\(1\) hypothesis on \(Y\) alone, so the explicit
    argument can be dropped.
\end{itemize}

\paragraph{Chapter structure.} \cref{sec:coheight_topology} proves the
topology-only lemma that coheight is preserved by open immersions
(\cref{lem:coheight_eq_of_isOpenEmbedding}). \cref{sec:coheight_spec_duality}
identifies coheight on a \(\Spec\) with the height of the underlying prime ideal
(\cref{lem:coheight_spec_eq_height_primeSpectrum}). \cref{sec:coheight_bridge}
assembles the bridge \cref{thm:ringKrullDim_stalk_eq_coheight}. The closing
\cref{sec:coheight_le_one_instance} packages the codim-\(1\) specialisation
\cref{lem:ringKrullDimLE_of_coheight_eq_one} that downstream files consume.

\section{Coheight is preserved by open immersions}
\label{sec:coheight_topology}

% Source pointer: \texttt{analogies/stacks-00tt-coheight.md} Decision 2 sketch
% L246--257 (the recipe extracts the lemma from a generic
% \texttt{Specializes.mem\_open} chain-lifting argument). The lemma is a
% generic topological fact about the specialisation preorder; no external
% reference shipped at \texttt{b80f227} states it under this exact
% formulation.

\begin{lemma}[Coheight on opens equals coheight in the ambient space]\leanok
  \label{lem:coheight_eq_of_isOpenEmbedding}
  \dcref{custom}
  \lean{Order.coheight_eq_of_isOpenEmbedding}
  Let \(X\) be a topological space carrying the specialisation preorder, and
  let \(U \subseteq X\) be an open subset. For any \(z \in U\),
  \[
    \operatorname{coheight}_{X}(z) \;=\; \operatorname{coheight}_{U}\bigl(\langle z, \, z \in U\rangle\bigr),
  \]
  where on the right \(U\) carries the subspace topology and \(\langle z, \,
  z \in U\rangle\) denotes \(z\) viewed as a point of \(U\).
\end{lemma}

\begin{proof}

\leanok
  Both sides are by definition the supremum, over natural numbers \(n\), of the
  set of \(n\) for which there exists a strictly increasing chain
  \(z = a_0 \prec a_1 \prec \cdots \prec a_n\) in the relevant
  specialisation preorder (with \(\prec\) denoting strict specialisation). It
  therefore suffices to give an order-preserving bijection between the
  strict-generisation chains starting at \(z\) in \(X\) and those starting at the
  corresponding point of \(U\).

  \emph{Chains in \(U\) embed into \(X\).} The inclusion \(U \hookrightarrow X\) is
  continuous, so it preserves specialisation: if \(y \rightsquigarrow x\) in \(U\)
  then \(y \rightsquigarrow x\) in \(X\), and distinct points of \(U\) remain
  distinct in \(X\). Hence every strict-generisation chain in \(U\) starting at
  \(z\) lifts to one in \(X\) of the same length.

  \emph{Chains in \(X\) starting at \(z\) all lie in \(U\).} Conversely, suppose
  \(z = a_0 \prec a_1 \prec \cdots \prec a_n\) is a strict-generisation chain
  in \(X\) with \(a_0 = z \in U\). For each \(i \geq 0\), the relation \(a_i \succeq
  z\) (in the specialisation preorder) is \(z \rightsquigarrow a_i\), i.e.\
  \(a_i\) is in the closure of \(\{z\}\) \emph{viewed from above}; in our
  convention the strict order \(a \prec b\) records \(b \rightsquigarrow a\)
  (the generic point is the strict maximum), so \(a_i \succeq z\) in \(X\) means
  \(a_i \rightsquigarrow z\) in \(X\). By the Mathlib lemma
  \texttt{Topology.Inseparable.Specializes.mem\_open} (any specialisation
  \(a \rightsquigarrow z\) with \(z\) in an open set \(U\) forces \(a \in U\), since
  \(U\) is open and contains \(z\), hence contains every point whose closure
  contains \(z\) --- equivalently every \(a\) with \(a \rightsquigarrow z\)),
  every \(a_i\) lies in \(U\). Thus the chain descends to a strict-generisation
  chain \(\langle z\rangle = \langle a_0\rangle \prec \langle a_1\rangle \prec
  \cdots \prec \langle a_n\rangle\) in \(U\) of the same length.

  These two operations are mutually inverse on the set of strict-generisation
  chains starting at \(z\) (resp.\ at \(\langle z, \, z \in U\rangle\)) and
  preserve length, so the suprema agree.
\end{proof}

\section{Coheight on Spec equals height in PrimeSpectrum}
\label{sec:coheight_spec_duality}

% Source pointer: \texttt{analogies/stacks-00tt-coheight.md} Decision 2 sketch
% L259--268. The convention reconciliation step traces to Mathlib's
% \texttt{spec\_le\_iff} in \texttt{AlgebraicGeometry.AffineSpace} L438--447
% and to \texttt{Order.coheight\_orderIso} +
% \texttt{Order.height\_toDual} in \texttt{Mathlib.Order.KrullDimension}.

\begin{lemma}[Coheight on \(\Spec R\) equals height in \(\operatorname{PrimeSpectrum} R\)]\leanok
  \label{lem:coheight_spec_eq_height_primeSpectrum}
  \dcref{custom}
  \lean{Order.coheight_spec_eq_height_primeSpectrum}
  Let \(R\) be a commutative ring and let \(p \in \Spec R\) (viewed as a point of
  the scheme spectrum, with the specialisation preorder). Then
  \[
    \operatorname{coheight}_{\Spec R}(p) \;=\; \operatorname{height}_{\operatorname{PrimeSpectrum}(R)}(p),
  \]
  where the right-hand side is the supremum of the lengths of strict
  prime-ideal chains \(\mathfrak{p}_0 \subsetneq \cdots \subsetneq \mathfrak{p}\)
  ending at the prime \(\mathfrak{p}\) corresponding to \(p\).
\end{lemma}

\begin{proof}

\leanok
  The specialisation preorder on \(\Spec R\) is the dual of the natural
  inclusion preorder on \(\operatorname{PrimeSpectrum}(R)\). Concretely, for
  \(p, q \in \Spec R\) with underlying primes \(\mathfrak{p}, \mathfrak{q}\),
  the relation \(p \leq q\) in the specialisation preorder (which is
  \(q \rightsquigarrow p\) in topological notation) corresponds to the prime
  containment \(\mathfrak{q} \subseteq \mathfrak{p}\). This is the content of
  the Mathlib lemma \texttt{spec\_le\_iff} (located in
  \texttt{Mathlib.AlgebraicGeometry.AffineSpace}, lines 438--447 at commit
  \texttt{b80f227}, with the explicit remark that the preorder on \(\Spec R\)
  equals the order-dual of the preorder on \(\operatorname{PrimeSpectrum} R\)).

  Under this order anti-isomorphism, coheight of \(p\) in \(\Spec R\)
  (= supremum of strict-generisation chain lengths starting at \(p\),
  i.e.\ ascending chains \(p \prec a_1 \prec \cdots \prec a_n\) in the
  specialisation preorder) translates to height of \(\mathfrak{p}\) in
  \(\operatorname{PrimeSpectrum}(R)\) (= supremum of strict descending prime-ideal
  chains \(\mathfrak{p} \supsetneq \mathfrak{q}_1 \supsetneq \cdots \supsetneq
  \mathfrak{q}_n\), equivalently strict ascending chains
  \(\mathfrak{q}_n \subsetneq \cdots \subsetneq \mathfrak{q}_1 \subsetneq
  \mathfrak{p}\)). The transport is mechanised by the Mathlib lemmas
  \texttt{Order.coheight\_orderIso} (coheight is invariant under order
  isomorphisms) and \texttt{Order.height\_toDual} (height in the order-dual
  equals coheight in the original order), applied to the order-dual
  identification \(\Spec R \simeq_o (\operatorname{PrimeSpectrum} R)^{\operatorname{op}}\)
  furnished by \texttt{spec\_le\_iff}.
\end{proof}

\section{The coheight \(=\) Krull-dim-of-stalk bridge}
\label{sec:coheight_bridge}

% Source pointer: \texttt{analogies/stacks-00tt-coheight.md} Decision 2 sketch
% L280--296 (the 5-step assembly recipe). Mathlib API surface consumed:
%   - \texttt{AlgebraicGeometry.exists\_isAffineOpen\_mem\_and\_subset}
%     (\texttt{AffineScheme.lean:271})
%   - \texttt{IsAffineOpen.primeIdealOf}, \texttt{IsAffineOpen.isLocalization\_stalk}
%     (\texttt{AffineScheme.lean:806--813})
%   - \texttt{IsLocalization.AtPrime.ringKrullDim\_eq\_height}
%     (\texttt{Mathlib.RingTheory.Ideal.Height:341--348})
%   - \texttt{Ideal.height\_eq\_primeHeight}
%     (\texttt{Mathlib.RingTheory.Ideal.Height:45})
%   - \cref{lem:coheight_eq_of_isOpenEmbedding}, \cref{lem:coheight_spec_eq_height_primeSpectrum}

\begin{theorem}[Coheight-to-Krull-dim bridge for a scheme point]\leanok
  \label{thm:ringKrullDim_stalk_eq_coheight}
  \dcref{custom}
  \lean{AlgebraicGeometry.Scheme.ringKrullDim_stalk_eq_coheight}
  \uses{lem:coheight_eq_of_isOpenEmbedding, lem:coheight_spec_eq_height_primeSpectrum}

  Let \(X\) be a scheme and let \(z \in X\) be any point. Then
  \[
    \operatorname{ringKrullDim}(\mathcal{O}_{X,z}) \;=\; \operatorname{coheight}_{X}(z),
  \]
  i.e.\ the Krull dimension of the stalk at \(z\) equals the coheight of \(z\) in
  the underlying topological space of \(X\) (with respect to the specialisation
  preorder).
\end{theorem}

\begin{proof}

  \leanok
  The proof proceeds in five steps, assembling Mathlib API rather than
  arguing from a single textbook reference.

  \emph{Step 1: pick an affine open neighbourhood of \(z\).} By
  \texttt{AlgebraicGeometry.exists\_isAffineOpen\_mem\_and\_subset}
  applied to the open set \(X\) itself, there exists an affine open
  \(U \subseteq X\) with \(z \in U\). Write \(\langle z, \, z \in U \rangle \in U\)
  for \(z\) viewed as a point of the subscheme \(U\).

  \emph{Step 2: name the prime corresponding to \(z\) inside \(U\).} Since
  \(U \cong \Spec \Gamma(U, \mathcal{O}_X)\) via the affine equivalence
  \(U \simeq \Spec R\) (with \(R := \Gamma(U, \mathcal{O}_X)\)), the point
  \(\langle z, \, z \in U\rangle\) corresponds to a prime ideal \(\mathfrak{p}
  \subseteq R\). In Mathlib this prime is named
  \texttt{IsAffineOpen.primeIdealOf}\,\(\langle z, \, z \in U\rangle\).

  \emph{Step 3: the stalk is the localisation of \(\Gamma(U)\) at
  \(\mathfrak{p}\).} The classical affine-stalk identification says
  \(\mathcal{O}_{X,z} \cong R_{\mathfrak{p}}\). In Mathlib this is the
  instance \texttt{IsAffineOpen.isLocalization\_stalk}, which witnesses
  \(\mathcal{O}_{X,z}\) as a localisation of \(R\) at \(\mathfrak{p}\) in the
  \texttt{IsLocalization.AtPrime} sense.

  \emph{Step 4: Krull dim of the localisation equals the height of
  \(\mathfrak{p}\).} The Mathlib lemma
  \texttt{IsLocalization.AtPrime.ringKrullDim\_eq\_height} gives
  \[
    \operatorname{ringKrullDim}(\mathcal{O}_{X,z}) \;=\;
    \operatorname{ringKrullDim}(R_{\mathfrak{p}}) \;=\;
    \operatorname{height}(\mathfrak{p}),
  \]
  where \(\operatorname{height}(\mathfrak{p})\) is the prime-ideal height in
  \(\operatorname{PrimeSpectrum}(R)\). By
  \texttt{Ideal.height\_eq\_primeHeight} this matches the
  \(\operatorname{height}_{\operatorname{PrimeSpectrum}(R)}\) used in
  \cref{lem:coheight_spec_eq_height_primeSpectrum}.

  \emph{Step 5: lift back from \(U\) to \(X\).} By
  \cref{lem:coheight_spec_eq_height_primeSpectrum} applied to \(R\) and the
  affine identification \(U \simeq \Spec R\),
  \[
    \operatorname{height}_{\operatorname{PrimeSpectrum}(R)}(\mathfrak{p}) \;=\;
    \operatorname{coheight}_{\Spec R}(\langle\text{image of } z\rangle) \;=\;
    \operatorname{coheight}_{U}\bigl(\langle z, \, z \in U\rangle\bigr).
  \]
  Then \cref{lem:coheight_eq_of_isOpenEmbedding} applied to the open
  \(U \subseteq X\) promotes the right-hand side back to
  \(\operatorname{coheight}_{X}(z)\). Concatenating the equalities from
  Steps~4--5 gives the conclusion.
\end{proof}

\section{Codim-1 specialisation: synthesising \texorpdfstring{\(\operatorname{KrullDimLE} 1\)}{KrullDimLE 1}}
\label{sec:coheight_le_one_instance}

% Source pointer: \texttt{analogies/stacks-00tt-coheight.md} Decision 2
% recipe scaffold L298--307 + Decision 2 hreg\_dim refactor sketch
% L319--340. This block is the consumer-facing wrapper of the bridge
% \cref{thm:ringKrullDim_stalk_eq_coheight} for the codim-\(1\) case used in
% \cref{chap:Albanese_CodimOneExtension} and
% \cref{chap:RiemannRoch_WeilDivisor}.

\begin{lemma}[Coheight \(= 1\) forces stalk \(\operatorname{Krull\,dim} \leq 1\)]\leanok
  \label{lem:ringKrullDimLE_of_coheight_eq_one}
  \dcref{custom}
  \lean{AlgebraicGeometry.Scheme.ringKrullDimLE_of_coheight_eq_one}
  \uses{thm:ringKrullDim_stalk_eq_coheight}

  Let \(X\) be a scheme and let \(z \in X\) be a point with
  \(\operatorname{coheight}_{X}(z) = 1\). Then the stalk \(\mathcal{O}_{X,z}\)
  has Krull dimension at most \(1\), equivalently \(\operatorname{Ring.KrullDimLE}\,1\,
  \mathcal{O}_{X,z}\) holds.
\end{lemma}

\begin{proof}

  \leanok
  By \cref{thm:ringKrullDim_stalk_eq_coheight},
  \(\operatorname{ringKrullDim}(\mathcal{O}_{X,z}) = \operatorname{coheight}_{X}(z) = 1\).
  The Mathlib predicate \(\operatorname{Ring.KrullDimLE}\,1\,R\) unfolds to
  \(\operatorname{ringKrullDim}(R) \leq 1\), so the inequality \(1 \leq 1\)
  discharges the goal.
\end{proof}

\section{Consumer wiring and the Stacks 00TT carve-out}
\label{sec:coheight_consumers}

\paragraph{Consumer 1: the \texttt{hreg\_dim} refactor in
\texorpdfstring{\cref{chap:Albanese_CodimOneExtension}}{the codim-1 extension chapter}.}
The private helper
\texttt{smooth\_codim\_one\_maximalIdeal\_isPrincipal\_and\_ne\_bot}
inside that chapter's Lean target currently carries an internal
\texttt{sorry} of the form
\[
  \text{hreg\_dim : IsRegularLocalRing}\,\mathcal{O}_{X,z}\;\wedge\;
  \operatorname{ringKrullDim}\,\mathcal{O}_{X,z} = 1.
\]
After this chapter's
\cref{lem:ringKrullDimLE_of_coheight_eq_one} (and its underlying
\cref{thm:ringKrullDim_stalk_eq_coheight}) land, the conjunction splits into
two independent goals:
\begin{itemize}
  \item The right conjunct \(\operatorname{ringKrullDim}(\mathcal{O}_{X,z}) = 1\)
    is discharged by
    \(\operatorname{coheight}_{X}(z) = 1\) (the standing codim-\(1\) hypothesis
    on \(z\)) plus \cref{thm:ringKrullDim_stalk_eq_coheight}; the antisymmetric
    pair \(\operatorname{ringKrullDim} \leq 1\) (\cref{lem:ringKrullDimLE_of_coheight_eq_one})
    and \(\operatorname{ringKrullDim} \geq 1\) (also from the bridge, since
    \(1 \leq 1\)) gives equality.
  \item The left conjunct \(\operatorname{IsRegularLocalRing}\,\mathcal{O}_{X,z}\)
    is the genuinely-irreducible \emph{Stacks tag \texttt{00TT}} statement
    ``smooth ring map over a field forces regular stalks''. It is a separate,
    out-of-scope sub-project (Mathlib does \emph{not} ship the implication
    \texttt{Algebra.Smooth k A \(\to\) IsRegularLocalRing \(A_{\mathfrak{p}}\)} at
    \texttt{b80f227}; the project-side build is \(\sim 200\)--\(300\)~LOC of
    cotangent-space and pure-dimension machinery, not packaged in this
    chapter).
\end{itemize}
Thus the present chapter \emph{halves} the \texttt{hreg\_dim} obligation in
\cref{chap:Albanese_CodimOneExtension}: one conjunct closes via the bridge,
one remains as a named Mathlib gap.

\paragraph{Consumer 2: the \texttt{Scheme.RationalMap.order} threading
hygiene in \texorpdfstring{\cref{chap:RiemannRoch_WeilDivisor}}{the Weil-divisor chapter}.}
The order-at-a-prime-divisor map currently carries an explicit instance
argument
\[
  [\,\operatorname{Ring.KrullDimLE}\,1\,(X.\text{presheaf.stalk}\,Y.\text{point})\,]
\]
at every use site. Because \(Y\) is a prime divisor of \(X\) (so
\(\operatorname{coheight}_{X}(Y.\text{point}) = 1\) by the definition of a prime
divisor as a codim-\(1\) irreducible closed subset),
\cref{lem:ringKrullDimLE_of_coheight_eq_one} synthesises the instance
automatically from the codim-\(1\) hypothesis on \(Y\), so the explicit instance
argument can be dropped from \texttt{Scheme.RationalMap.order}'s signature.

\paragraph{Consumer 3: future codim-\(1\) infrastructure.} Any subsequent
project lemma that needs to discharge ``the stalk at a codim-\(1\) point has
Krull dim \(\leq 1\)'' --- the implicit hypothesis behind every DVR-at-codim-\(1\)
argument (cf.\ \cref{lem:smooth_codim_one_dvr} of
\cref{chap:Albanese_CodimOneExtension}) --- now has a one-liner consequence.
The instance form
\cref{lem:ringKrullDimLE_of_coheight_eq_one} is the right shape for
typeclass-search to find it without ceremony.

\section{Lean encoding}
\label{sec:coheight_lean_encoding}

The Lean target file is the new
\texttt{AlgebraicJacobian/Albanese/CoheightBridge.lean}. The declarations to
introduce, in order:

\begin{enumerate}
  \item \texttt{Order.coheight\_eq\_of\_isOpenEmbedding} ---
    \cref{lem:coheight_eq_of_isOpenEmbedding}. Signature (informal): for
    \(X\) a topological space, \(U \subseteq X\) open, and \(z \in U\),
    \(\operatorname{coheight}_{X}(z) = \operatorname{coheight}_{U}(\langle z, hz\rangle)\).
    Tools: \texttt{Topology.Inseparable.Specializes.mem\_open} for the
    chain-lifting direction; \texttt{Subtype.val} continuity for the
    embedding direction. Estimated \(\sim 20\)--\(30\)~LOC.
  \item \texttt{Order.coheight\_spec\_eq\_height\_primeSpectrum} ---
    \cref{lem:coheight_spec_eq_height_primeSpectrum}. Signature (informal):
    for \(R\) a \texttt{CommRingCat} and \(p : \Spec R\) a point, coheight of \(p\)
    in \(\Spec R\) equals height in \(\operatorname{PrimeSpectrum} R\). Tools:
    \texttt{spec\_le\_iff} (\texttt{Mathlib.AlgebraicGeometry.AffineSpace} L438--447),
    \texttt{Order.coheight\_orderIso}, \texttt{Order.height\_toDual}
    (\texttt{Mathlib.Order.KrullDimension}). Estimated \(\sim 10\)--\(15\)~LOC.
  \item \texttt{AlgebraicGeometry.Scheme.ringKrullDim\_stalk\_eq\_coheight} ---
    \cref{thm:ringKrullDim_stalk_eq_coheight}. Signature (informal): for
    \(X : \texttt{Scheme}\) and \(z : X\),
    \(\operatorname{ringKrullDim}(X.\text{presheaf.stalk}\,z) =
    \operatorname{coheight}\,z\). Five-step assembly per the proof above.
    Tools: \texttt{exists\_isAffineOpen\_mem\_and\_subset},
    \texttt{IsAffineOpen.primeIdealOf}, \texttt{IsAffineOpen.isLocalization\_stalk},
    \texttt{IsLocalization.AtPrime.ringKrullDim\_eq\_height},
    \texttt{Ideal.height\_eq\_primeHeight}, plus the two lemmas above.
    Estimated \(\sim 30\)--\(50\)~LOC.
  \item \texttt{AlgebraicGeometry.Scheme.ringKrullDimLE\_of\_coheight\_eq\_one} ---
    \cref{lem:ringKrullDimLE_of_coheight_eq_one}. Signature (informal):
    given \(X : \texttt{Scheme}\), \(z : X\), and \(\operatorname{coheight}\,z = 1\),
    derive \texttt{Ring.KrullDimLE 1 (X.presheaf.stalk z)}. Trivial corollary
    of (3). Estimated \(\sim 10\)~LOC. Lemma vs.\ instance phrasing: the
    \texttt{coheight\,z = 1} hypothesis is not a typeclass, so this lands as
    a plain lemma. Downstream files that have the codim-\(1\) hypothesis in
    hand invoke it explicitly; downstream typeclass infrastructure that
    bundles ``\(z\) is a codim-\(1\) point'' into a typeclass (if any is added
    later) can re-export as an instance.
\end{enumerate}

\paragraph{Mathlib readiness audit.} At commit \texttt{b80f227} the prover
should verify:
\begin{itemize}
  \item \texttt{Order.coheight}, \texttt{Order.height}, \texttt{Order.coheight\_orderIso},
    \texttt{Order.height\_toDual} all live in
    \texttt{Mathlib.Order.KrullDimension} and are imported transitively by
    \texttt{Mathlib.AlgebraicGeometry.Stalk}.
  \item \texttt{Topology.Inseparable.Specializes.mem\_open} lives in
    \texttt{Mathlib.Topology.Inseparable}; the file already imports it via
    \texttt{Mathlib.Topology.Defs.Filter}.
  \item \texttt{IsAffineOpen.isLocalization\_stalk} and
    \texttt{IsAffineOpen.primeIdealOf} live in
    \texttt{Mathlib.AlgebraicGeometry.AffineScheme} (lines 806--813 and
    surrounding, at the pinned commit).
  \item \texttt{IsLocalization.AtPrime.ringKrullDim\_eq\_height} and
    \texttt{Ideal.height\_eq\_primeHeight} live in
    \texttt{Mathlib.RingTheory.Ideal.Height} (lines 341--348 and 45 at the
    pinned commit).
  \item \texttt{Ring.KrullDimLE} is defined and unfolds to
    \(\operatorname{ringKrullDim}(R) \leq n\) at \texttt{Mathlib.RingTheory.KrullDimension.Basic}.
\end{itemize}
No new core axiom is required. The chapter introduces no new macro.

\section{Out of scope}
\label{sec:coheight_out_of_scope}

This chapter packages \emph{only} the coheight-to-Krull-dim bridge and its
codim-\(1\) specialisation. The following adjacent content is explicitly out
of scope:
\begin{itemize}
  \item \textbf{Stacks tag \texttt{00TT} (smooth \(\Rightarrow\) regular stalks).}
    This is a separate sub-project of \(\sim 200\)--\(300\)~LOC, tracked in
    \texttt{analogies/stacks-00tt-coheight.md} Decision 1. The directive
    flags it as iter-200+ work. The natural project landing site, when
    attacked, is a standalone file
    \texttt{AlgebraicJacobian/Albanese/SmoothToRegular.lean} (not part of
    this chapter).
  \item \textbf{The codim-\(1\) valuative criterion / Milne Lemma 3.3
    difference-map construction.} This is the second multi-iter sub-project
    flagged in \texttt{analogies/stacks-00tt-coheight.md} Decision 3
    (\(\sim 300\)--\(500\)~LOC of project-side new infrastructure); see
    \cref{lem:milne_codim1_indeterminacy} in
    \cref{chap:Albanese_CodimOneExtension} for the chapter that owns it.
  \item \textbf{Refactoring \texttt{Scheme.RationalMap.order} to drop its
    explicit \texttt{[Ring.KrullDimLE 1 \_]} instance argument.} This is a
    downstream consequence of \cref{lem:ringKrullDimLE_of_coheight_eq_one}
    that lives in \cref{chap:RiemannRoch_WeilDivisor}; the present chapter
    only \emph{enables} the refactor, it does not perform it.
  \item \textbf{Mathlib upstream PRs.} The two lemmas
    \cref{lem:coheight_eq_of_isOpenEmbedding} and
    \cref{thm:ringKrullDim_stalk_eq_coheight} are natural candidates for
    upstreaming (their natural Mathlib homes are
    \texttt{Mathlib.Order.KrullDimension} and a new
    \texttt{Mathlib.AlgebraicGeometry.Coheight} respectively); the chapter
    flags this as an observation but does not pursue it.
  \item \textbf{General codim-\(n\) specialisations.} Only the codim-\(1\)
    instance is packaged here, because that is the consumer-facing case.
    A general \texttt{ringKrullDimLE\_of\_coheight\_le} or
    \texttt{ringKrullDimLE\_of\_coheight\_eq} family of lemmas can be added
    iteratively if downstream demand emerges, but is not required for any
    A.4.* sub-row at the time of writing.
\end{itemize}
